%%=============================================================================
%% Inleiding
%%=============================================================================

\chapter{\IfLanguageName{dutch}{Inleiding}{Introduction}}%
\label{ch:inleiding}
Deze bachelorproef onderzoekt Master Data Management (MDM), multi-ERP-omgevingen en het concept van "golden records''. Specifiek richt het onderzoek zich op de casus van IPCOM, een bedrijf in de isolatiesector dat met meerdere ERP-systemen werkt. IPCOM heeft op dit moment het probleem dat verschillende entiteiten binnen de groep (zoals in Oostenrijk, Zwitserland en Noorwegen) dezelfde artikelen van dezelfde leveranciers kopen, maar dat deze door de verschillende bronsystemen niet eenvoudig te identificeren zijn als hetzelfde product. Hierdoor ontstaan inconsistenties in artikel- en leveranciersdata, wat een directe impact heeft op rapportering en besluitvorming.

Om dit probleem aan te pakken is IPCOM momenteel een pilot gestart met de commerciële MDM-tool Profisee. Dit roept een interessante nieuwe vraag op voor deze bachelorproef. In plaats van een theoretische oplossing te bedenken, wordt er een proof-of-concept (PoC) gebouwd in Microsoft Fabric die als objectieve benchmark dient. Hoe presteert een eigen, in-house gebouwde oplossing op basis van deterministische regels en machine learning in vergelijking met de Profisee-setup, wanneer exact dezelfde inputdata en mappings worden gebruikt? Hoe past zo'n MDM-tool in een datawarehouse-architectuur en hoe kan de datakwaliteit over de verschillende systemen heen worden gemeten?

De belangrijkste doelgroep wordt gevormd door data-engineers en business-gebruikers binnen de organisatie die werken met deze data in het datawarehouse en het rapportageplatform (bijvoorbeeld Power BI). Zij ondervinden dagelijks hinder van gefragmenteerde data en hebben nood aan een gecentraliseerde, betrouwbare bron van waarheid, namelijk een "golden record" per artikel en leverancier, die de datakwaliteit structureel verbetert.

De centrale onderzoeksvraag luidt als volgt: Hoe kan een in-house Master Data Management-oplossing in Microsoft Fabric worden ontworpen en geïmplementeerd voor artikel- en leveranciersdata, zodat deze kan dienen als objectieve benchmark tegenover een commerciële MDM-tool (Profisee) binnen een multi-ERP-omgeving?
De doelstelling van deze bachelorproef is eenduidig: een PoC bouwen voor een MDM-tool in Microsoft Fabric, specifiek gericht op de ERP-data van Oostenrijk, Zwitserland en Noorwegen. Op basis van de vergelijking tussen deze PoC en de Profisee-pilot kunnen concrete aanbevelingen worden geformuleerd over de toegevoegde waarde van commerciële pakketten versus in-house oplossingen voor data harmonisatie en matching.

\subsection{Deelvragen m.b.t probleemdomein}
\begin{itemize}
    \item Hoeveel duplicaten en conflicterende records komen voor over de geselecteerde ERP-systemen heen?
    \item Welke datakwaliteitsproblemen treden het meest op bij de geselecteerde artikel- en leveranciersdata (bijvoorbeeld het hoge percentage ontbrekende waarden)?
    \item Welke identificatievelden (zoals EAN en artikelnaam) zijn bruikbaar in de systemen om een correcte mapping op te zetten?
    \item Hoe vaak verandert een golden record de oorspronkelijke data bij welk specifiek bronsysteem?
\end{itemize}

\subsection{Deelvragen m.b.t oplossingsdomein}
\begin{itemize}
    \item Hoe presteert de in-house Fabric PoC ten opzichte van de commerciële Profisee-setup wanneer dezelfde mappings en rulesets worden gebruikt?
    \item Welke matching-technieken (deterministisch op basis van EAN versus fuzzy matching of AI op artikelnaam) leveren de hoogste nauwkeurigheid?
    \item Welke consolidatieregels (survivorship-logica) zorgen voor de meest consistente "golden records"?
    \item Hoe kan de publicatie van golden records worden ingericht zodat lokale id's correct verwijzen naar de golden records en alle entiteiten hiervan profiteren?
    \item Hoeveel verbetering in datakwaliteit wordt gerealiseerd en visueel meetbaar gemaakt na implementatie van de MDM-oplossing?
\end{itemize}

\section{\IfLanguageName{dutch}{Opzet van deze bachelorproef}{Structure of this bachelor thesis}}%
\label{sec:opzet-bachelorproef}

% Het is gebruikelijk aan het einde van de inleiding een overzicht te
% geven van de opbouw van de rest van de tekst. Deze sectie bevat al een aanzet
% die je kan aanvullen/aanpassen in functie van je eigen tekst.

De rest van deze bachelorproef is als volgt opgebouwd:

In Hoofdstuk \ref{ch:stand-van-zaken} wordt een overzicht gegeven van de theoretische kaders rond Master Data Management, golden records en datakwaliteit binnen een multi-ERP-context.

In Hoofdstuk \ref{ch:methodologie} wordt de methodologie toegelicht voor het ontwikkelen van een proof-of-concept in Microsoft Fabric, waarbij artikel- en leveranciersdata uit drie ERP-systemen worden geharmoniseerd

In Hoofdstuk \ref{ch:conclusie} worden de resultaten van de proof-of-concept geëvalueerd en afgezet tegen een commerciële benchmark om een onderbouwd "build vs. buy"-advies te formuleren. 